\section{Introduction}
\label{sec:introduction}

Unlabeled data are often abundant, whereas obtaining labels may
require costly manual annotation. Self-training addresses this issue
by using a model's predictions as pseudo-labels for unlabeled
observations, which then contribute to subsequent training
\cite{lee2013pseudolabel}.

We write $\mathbb N=\{0,1,2,\ldots\}$,
$[\![a,b]\!]=\{a,a+1,\ldots,b\}$ for integers $a\leq b$,
and $\mathscr L(X)$ for the law of a random element $X$.

Given partially labeled observations
$\mathcal D:=\{(X_i,\Delta_iY_i):i\in[\![1,n]\!]\}$, where
$X_i\in\mathbb R^d$, $Y_i\in\{-1,+1\}$, and
$\Delta_i\in\{0,1\}$ indicates whether the label is observed,
we study a self-training algorithm that produces affine classifiers
of the form
\begin{equation}
    f_\theta(x)
    :=
    \operatorname{sign}\left(
        b+\frac{x^\top w}{\sqrt d}
    \right),
    \qquad
    \theta=(b,w)\in\mathbb R\times\mathbb R^d,
    \label{eq:def-classifier}
\end{equation}
with $\operatorname{sign}(0)=+1$.
Let
$I_{\mathcal U}:=\{i\in[\![1,n]\!]:\Delta_i=0\}$
denote the indices of the unlabeled observations.
At each iteration $t$, the current classifier assigns to each
$j\in I_{\mathcal U}$ a training score and a pseudo-label
\[
    r_j^t
    :=
    b^t+\frac{X_j^\top w^t}{\sqrt d},
    \qquad
    \widehat Y_j^t
    :=
    \operatorname{sign}(r_j^t).
\]
Given confidence thresholds $\kappa_-<0<\kappa_+$,
only observations with sufficiently confident scores are retained,
forming the pseudo-labeled dataset
\begin{align}
    \mathcal P_t
    :=
    \bigl\{(X_j,\widehat Y_j^t):j\in I_t\bigr\},
    \qquad
    I_t
    :=
    \left\{
        j\in I_{\mathcal U}:
        r_j^t\le\kappa_-
        \ \text{or}\
        r_j^t\ge\kappa_+
    \right\}.
    \label{eq:pl-dataset}
\end{align}
The model is then updated by taking one gradient step,
\begin{align}
    \theta^{t+1}
    :=
    \theta^t
    -
    \eta P_d
    \nabla_\theta L(\theta^t;\mathcal D_L,\mathcal P_t), 
    \label{eq:pl-gd-step}
\end{align}
where
$\mathcal D_L:=\{(X_i,Y_i):\Delta_i=1\}$
is the labeled dataset, $L$ is the surrogate objective
specified in \eqref{eq:ST-objective}, and $P_d$ is a
preconditioning matrix. Its pseudo-labeled contribution is
weighted by $\pi$ and normalized by the full unlabeled count $M$.
This procedure is repeated, with the same labeled and unlabeled
observations reused throughout training.

Unlike ordinary empirical risk minimization, the training objective
changes at each iteration, as the pseudo-labels and selected
observations depend on the current classifier. These adaptive
changes introduce discontinuities into the update rule and
nonconvexity into the underlying optimization problem.
The statistical analysis presents additional difficulties.
Because the same random observations are reused throughout
training, successive parameter updates are coupled through
their dependence on the data and previous iterates.
Furthermore, we are interested in understanding these dynamics
in the high-dimensional regime, where many tools from classical
asymptotic statistics are no longer applicable.

These considerations raise two natural questions:
\begin{itemize}
    \item Does the self-training algorithm converge, and what
    can be said about its limiting point?
    \item Can the dynamics of the parameter iterates and their
    statistical performance be characterized precisely in
    high dimensions?
\end{itemize}

\subsection{Contributions}

We address these questions by establishing convergence of the
self-training algorithm and deriving an asymptotically exact
characterization of its finite-time dynamics under a Gaussian-mixture
model in the high-dimensional regime.
The precise setting and assumptions are presented in
Section~\ref{sec:setting}. Our main contributions are threefold:

\begin{enumerate}
    \item \emph{Finite-dimensional convergence.}
    Under suitable assumptions, we prove that the self-training
    parameters converge for a fixed dataset, with an eventual
    geometric rate of convergence. We further characterize the
    limiting point in terms of the final pseudo-labeling and
    selection configuration.

    \item \emph{Asymptotically exact dynamics.}
    In the proportional asymptotic regime, we
    characterize the limiting empirical laws of
    parameter-coordinate trajectories and sample-score
    trajectories over every fixed time horizon through
    two effective stochastic processes.

    \item \emph{Long-time convergence.}
    Under additional assumptions, we establish convergence of
    the effective parameter and score processes as the number
    of iterations tends to infinity, with the high-dimensional
    limit taken first. For a representative unlabeled observation,
    the selection state and the assigned label when selected
    eventually stabilize almost surely.
\end{enumerate}

Our results provide foundations for further analysis of
self-training algorithms.
Section~\ref{sec:numerical} takes initial steps toward
understanding improvements in generalization and the effects
of class imbalance.
A comprehensive quantitative analysis of these mechanisms,
and of the trade-off between pseudo-label reliability and
the fraction of selected observations, remains open.

Section~\ref{sec:related-work} reviews self-training,
self-distillation, adaptive training procedures, and
high-dimensional dynamics.
Section~\ref{sec:setting} specifies the model and algorithm.
Section~\ref{sec:main-result} presents the finite-dimensional,
mean-field, and long-time results, proved respectively in
Appendices~\ref{sec:proof-finite-dim},
\ref{sec:proof-mean-field}, and \ref{sec:proof-long-time}.
Section~\ref{sec:numerical} presents numerical validation and
statistical interpretation, with proofs in
Appendix~\ref{sec:proof-numerical}.
Appendix~\ref{sec:math-prelim} collects mathematical preliminaries 
that might be useful to the reader.
